\documentclass[runningheads,draft]{llncs}
\usepackage{graphicx}
\usepackage{todonotes}
\usepackage{hyperref}
\newcommand{\till}[1]{\textcolor{red}{#1}}
\begin{document}
\title{A Test Paper}
\author{Jane Doe}
\institute{Somewhere}
\maketitle
\begin{abstract}
We study Autonomous Driving (AD) with Large Language Models (LLMs).
\end{abstract}

\section{Introduction}\label{sec:intro}
Autonomous vehicles are complex systems \cite{smith2020}.
Scenario-based Testing (SBT) is a common approach \cite{doe2021}\cite{smith2020}.
As shown in Figure~\ref{fig:arch}, the AD stack is large.
Fig.~\ref{fig:arch} shows the architecture.
In Section \ref{sec:method} we describe the method.
We use LLMs here. The behaviour is modelled in \cite{doe2021}.
This is a real-world problem and the real world is hard; see the real world data.
It's a lot of work, e.g. for 3 sensors running at 10ms.
The results are shown in pages 3-5.\cite{doe2021}
\cite{smith2020} shows that "quotes" matter.
We use Scenario-based Testing again. The color of the colour.
Tab.~\ref{tab:res} lists results. The optimized and optimised variants.
\till{rewrite this} TODO check this. See https://example.org/foo for details.
The OpenAI api and the Openai API.\footnote{\url{https://github.com/foo}}
Another note.\footnote{\url{https://github.com/foo}}

The remainder of this paper is structured as follows: Sect.~\ref{sec:method} presents the method, Sect.~\ref{sec:related} related work, and Sect.~\ref{sec:intro} again.

\section{Related Work}\label{sec:related}
Related work is related \cite{missingkey} and \cite{}.

\section{Method}\label{sec:method}
Our method is described here in detail with many words so that the section is filled
and does not trigger the short section hint for the method section of the paper at all.
It has more text. Scenario-based Testing (SBT) is defined again here.
\begin{figure}[t]
\includegraphics{arch}
\caption{Architecture.}\label{fig:arch}
\end{figure}
\begin{table}[t]
\caption{Results.}\label{tab:res}
\end{table}
\begin{figure}[t]
\caption{Unused.}\label{fig:unused}
\end{figure}
See Figures~\ref{fig:arch} and~\ref{fig:unused} and Eq.~\ref{eq:x}.
Known approaches~\cite{doe2021,smith2020} differ.
Approaches of the same year~\cite{lee2021,doe2021} keep their order.
One entry has no year~\cite{smith2020,noyear}.


\section{Evaluation}\label{sec:eval}

\section{Conclusion}\label{sec:conc}
We conclude. Unsupervised learning is supervised. The analyses are done.
% This is a longer commented out block of text that the author
% forgot to remove before submission and it contains many words
% which should be detected by the leftovers check

\bibliographystyle{splncs04}
\bibliography{refs}
\end{document}
